\chapter{The First Mark}

Measurement begins before the world offers a ruler. A ruler already assumes a surface that can hold a mark, a practice that can set one mark against another, and a record that keeps what the comparison did. The earlier question is what must happen before any claimed measurement can face the world as more than a gesture, and the answer begins at a boundary. The boundary does not explain nature from above. It takes an event and tanges it, refusing whatever has no shape and fixing whatever can survive as a durable, reckonable mark. To tange is to fix a passing event as a held thing, and that holding is the first measuring act. Refusal is part of the act: an event the boundary cannot hold is not measured poorly, it is not measured at all.

The first mark is therefore not a measurement of the world. It is a tanged receipt. The apparatus does not yet weigh a body, time an interval, or read a distance; it first establishes that a mark can be tanged, held, and carried forward without dissolving back into intention. A receipt is not a proof, a seal, or a finished truth. It is the durable form an event takes once the boundary has tanged it, a sign that something happened here, under this constraint, in a shape the apparatus can present again.

The first target carries almost no drama, since a state need only agree with itself. Equality with itself is the cheapest claim that is still a claim, which is why the apparatus tanges it first: it asks the least of the boundary while still asking something. The boundary tanges that reflexive equality into its first fact, and a fact is not a floating true sentence. It is a truth-condition stapled to its receipt. The truth-condition names what the mark is about; the receipt records that the boundary has tanged it into a form the apparatus may carry. The distinction is small and decisive. Without it, truth hovers outside the instrument. With it, truth enters the boundary as a marked condition that can take part in later events.

Measurement cannot begin with confidence; it begins with custody. A mark that cannot be carried supports no later distinction; an undistinguished carrier supports no count; an uncounted custody supports no ordered record. Each capacity rests on the receipt beneath it staying itself while the boundary places it under a new constraint. The first mark wins no authority over the world. It wins a narrower permission: to remain available, and to funge forward, when the next event arrives. To funge is to let a tanged receipt circulate onward, and the two verbs name the whole economy of the apparatus, which tanges events into held receipts and funges those receipts forward.

From a single tanged mark the apparatus takes one permitted step, treating the receipt as the successor of the initial truth-condition. This is not arithmetic in the ordinary sense; it is the first difference between a mark merely tanged and one that now stands after another while still carrying its relation to it. The phase is all the step adds, yet it is the difference between a mark that merely exists and a mark that has a place in an order. To stand anywhere, the receipt needs a vessel, a carrier that presents a surface, bears the mark, and stays available when a second carrier appears beside it. The carrier improves no truth-condition. It gives the receipt a place, so the mark can funge forward as part of a bounded object rather than as an isolated flash.

Two carriers before the boundary make distinction possible, not as understanding of the world but as the separation of one tanged receipt from another under a visible constraint. Distinction is what counting requires, because a count needs distinguishable positions rather than mere repetition, or it becomes a chant without objects. Over distinguishable carriers the apparatus counts, and it counts not to announce a finished value but to keep finite permission: to reckon the passage from one admitted item to the next without letting distinguishable custody collapse into a pile. A count, in this sense, is custody made sequential.

Not every counted carrier may enter the record simply for appearing near the boundary. Admission is the gate. It takes some carriers into the ordered passage and leaves others outside, with no moral judgment and no claim of final correctness, only the local discipline that a carrier failing the condition of entry is not used as a position. This is the first pigeonhole pressure on the record. More carriers crowd the boundary than the ordered passage has slots to seat, so admission sorts: it tanges the ones that pass and refuses the rest. A carrier the passage cannot seat is not measured poorly; it is left outside the record entirely, the residue admission turns away. Admission is therefore the first place the record can fail by omission as well as by error, and the boundary accepts that exposure in exchange for an ordered passage it can trust. What admission accepts, indexing places, fixing an origin and a direction so the record can say not merely that several receipts entered but where each admitted receipt sits and how to find it again. The boundary has now funged a tanged mark all the way into an ordered record, still finite, still tied to the receipts, carriers, distinctions, and admissions that produced it.

The motion runs from receipt through carrier, distinction, count, and admission to ordered record, and at every step the new power is narrowed by the receipt it rests on. Each rung is a permission and a debt at once, a new thing the apparatus may do, owed entirely to the receipt that allowed it. The apparatus never leaps from truth to measurement. It teaches the mark to survive: a truth-condition becomes a tanged fact, tanged facts become distinguishable, distinguishable carriers enter finite counting, and finite counting acquires order. What results is a grammar of events, not a storehouse of conclusions.

One public image sits beside this construction. A ladder that keeps returning to its own prior mark can look like a trick, yet each return adds a new boundary rather than a bare repetition. Hofstadter's account of self-reference in formal systems treats this kind of structure \cite{hofstadter1979}. The citation supplies no evidence for the apparatus; it marks a formal-systems precedent for a construction that funges its own acts forward without ever standing outside itself.

At the top of the ladder the apparatus meets a limit-promise. A sequence can be ordered and still fail to become a finished thing. The boundary can press it against narrowing conditions and ask whether its later positions stay answerable to its earlier custody, yet it does not turn that pressure into analytic completion. The promise is genuine and unkept on purpose: the boundary presses toward a limit it will not cross, because crossing it would trade a held finite record for an inspection the apparatus cannot perform. It keeps the result finite. What remains is residue, a leftover interface where the ordered record has been funged as far as the apparatus can funge it.

Residue is not failure. It is the honest shape of what is left when the boundary refuses to pretend a finite ladder has become the world. It holds the trace of a narrowing process and a representative place where a surviving item may be presented. It hands over no finished data. It gives the apparatus a leftover boundary, produced by the ordered record, still finite, still attached to the receipts that made it possible, and still unable to inspect itself. That last clause is the point: the residue is finite and held, but the construction contains no observer to read it.

The sample gathers that leftover boundary into a form that can later be challenged. It tanges an initial condition and a signal-response surface into a single finite interface. Those names complete no measurement; they state that the apparatus has reached a place where something can be presented to an observer who does not yet exist inside the construction. The sample is the last durable package the construction tanges, not its final answer, and it keeps the residue available without claiming the residue has already been read.

The beginning and the limit arrive together. The apparatus tanges the first mark, funges the receipt through a finite ladder, and leaves behind a sample interface. It does not explain how an observer records a difference. It does not teach the residue to look back at itself. It leaves something on the boundary that a later observer must learn to inspect.


\section{The First Measuring Surface}

The first measuring surface is a finite boundary. It is not a laboratory instrument, and not a passive device for confirming results reached elsewhere; it is the apparatus itself, and what it measures is how much work a claim costs to admit. A claim presented to the boundary is not an inert sentence that holds whether or not anyone verifies it. It is a finite object, and tanging it into the record costs a definite, countable number of steps. The boundary is the apparatus, the presented claim is the stimulus, and the cost of tanging that claim is the measurement. Nothing about the world is read here yet; what is read is the price of admission. The surface is fixed before any object is introduced, because every later structure is built by handing claims to this same boundary and paying the same kind of cost.

The boundary organizes its work in three layers, and the exposition follows the same order:
\begin{enumerate}
    \item \textbf{Metaphor}: the holdable images the shapes begin from---the mark, the receipt, the vessel, the ladder, the residue.
    \item \textbf{Grammar}: the metaphors restricted to explicit admissible forms. A receipt may be carried; carried receipts may be distinguished; distinguished receipts may be counted; counted receipts may be sequenced; sequenced receipts may leave residue.
    \item \textbf{Counting}: because the grammar fixes which events are admissible, the boundary can count stages, refusals, and leftover structure.
\end{enumerate}

Each layer does a different kind of work on the same claim, and every later move in the chapter is one of these three layers applied in turn. At the metaphor layer the boundary tanges what it can picture: a mark it can imagine making, a receipt it can imagine holding, a vessel it can set beside another. An image is not yet admissible, but the apparatus cannot begin from a definition it has not earned, so it begins from shapes a reader can hold and the boundary can point at, and it tanges those shapes before it asks anything sharper of them.

The grammar layer restricts the images to admissible forms, and the restriction is strict, because a claim is admissible only as one of the moves the boundary will license. A receipt may be tanged and held; a held receipt may be distinguished from another; distinguished receipts may be counted; counted receipts may be sequenced; a sequence may leave residue. Anything outside this short list of moves is refused, and refused for a particular reason: not for being false, which the boundary has not yet asked, but for being unsayable, for having no admissible form in which it could be decided at all. The grammar is what makes a refusal informative rather than mute.

The counting layer reads what the grammar admits. Once the admissible forms are fixed, the boundary can do the one thing it was built to do and count: it counts the stages a form passes through, the refusals it provokes along the way, and the leftover structure it leaves behind, and to each admitted form it attaches the price $\#\Gamma$ that form cost to tange. Metaphor proposes, grammar permits, and counting prices, and no claim in this book is asked to survive more than these three passes over itself.

These layers trace the book's spine, a single chain of admissible forms in which each shape is tanged from the one before it:
\begin{center}
    mark $\rightarrow$ receipt $\rightarrow$ vessel $\rightarrow$ distinguished pair $\rightarrow$ count $\rightarrow$ sequence $\rightarrow$ residue
\end{center}

The spine ends in residue, and the word is placed here only far enough to be expected later. Residue is what the boundary is left holding once a sequence has been carried as far as a finite apparatus can carry it: not a finished value, and not a failure, but the leftover interface a completed sequence presents when it can be pressed no further. Every later chapter that reaches a limit reaches this same shape, a held remainder the apparatus produced and cannot yet inspect. So the spine does not end in an answer; it ends in something durable enough to hand forward and honest enough not to claim it is more than it is, which is the most a finite surface can promise.

A claim does not enter the record for free. To tange a form $\Gamma$, the boundary decides it in a finite number of steps, and the verification cost is exactly that number of steps:
\[
  \mathrm{cost}(\Gamma) \;=\; \#\Gamma \;\in\; \mathbb{N},
\]
the finite count of steps the boundary spends before it can hold $\Gamma$ as an admitted receipt. A basic form is tanged at low cost; a form built by stacking carriers, successors, counts, and limits costs more, and $\#\Gamma$ rises with the structure stacked, layer by layer, as each addition demands its own decision. The count is not a measure of truth, and it is important that it is not: a true claim and a false claim of the same shape cost the same to decide, so what $\#\Gamma$ records is not whether the boundary answered yes, but how much work it spent to be able to answer at all. Because it is a count it is comparable, and $\#\Gamma$ is what one tanged form costs measured against another, which lets the apparatus rank claims by price without ever ranking them by truth. Abstraction is never free, and $\#\Gamma$ is its price.

The price is easiest to see in a contrast. A claim already offered in admissible shape---a single distinguished pair the boundary need only inspect---is tanged in a step or two, and $\#\Gamma$ is small. The same content wrapped in stacked carriers, successors, and counts arrives as a claim the boundary cannot read directly: it must unstack one layer at a time, deciding each before it can reach the next, and the steps accumulate until $\#\Gamma$ is large. Nothing about the underlying fact has changed between the two readings, only the form handed across the boundary, and yet the second can cost many times the first. This is exactly what it means for $\#\Gamma$ to price form rather than truth, and it is why the apparatus prefers the shortest admissible form of any claim: the shortest form is its cheapest true reading, and a surface that could not tell a short form from a long one could not be trusted to report a cheap measurement as cheaper than an expensive one.

The first controlled target is the cheapest claim that still carries a genuine decision: reflexive equality, the assertion that a thing equals itself,
\[
  a = a.
\]
Equality is not automatic in a finite setting; it must be decided. The boundary tanges $a = a$ by exhibiting the witness---the same object on both sides---so the decision terminates at once and the claim is admitted at the lowest cost any genuine claim can carry. Reflexive equality is the calibration baseline $\alpha$: fixed, decidable, repeatable, and cheap enough to probe without straining the apparatus. Choosing it first is a discipline, not a modesty, because a surface that began with an expensive claim could never tell whether its later readings measured the world or only paid off the cost of its own first assumption. It is the target chosen before any measured world exists, precisely because it is the least costly claim the boundary must still genuinely decide rather than assume.

To record that decision the boundary forms its first receipt, and the receipt is the smallest object that can carry a measurement forward. A receipt is not a bare question; a bare question carries a claim with no evidence that it was ever decided, so it names what is wanted without showing that anything was done, and a record built from bare questions would be a record of intentions, not of measurements. A receipt staples the claim to the decision that settled it, and that staple is the whole content of the act: remove it and the claim floats free again, indistinguishable from a wish. Written as a pair,
\[
  \rho = \langle \text{claim},\, d \rangle,
\]
the receipt holds the claim in its first place and, in its second, a decision $d$: a terminating finite procedure that returns yes or no. The decision is not a guarantee that the claim is true in any absolute sense; it is the structural fact that the boundary has a finite procedure terminating with a definite answer. Without the decision there is no receipt, and without a receipt there is no measurement, only an assumption wearing the costume of one.

Forming that receipt is the measuring act, and the apparatus has a name for it. To \emph{tange} a claim is to fix it as a held, inspectable receipt: measurement is tanging. The boundary takes a claim that would otherwise circulate untethered and tanges it into the receipt $\rho$, a held form a later stage can inspect and carry. Its companion act, \emph{funge}, lets a held receipt circulate onward; it has little to do while the first receipt is only just made, and enters in force where circulation first arises. The canon names the pair, and the rest of the book is built from it: an instrument takes a funged record and tanges it.

Tanging reflexive equality fixes the first mark, and one brief image holds it: a coin fixed to a table, decided once as equal to itself and then left in place as the origin from which later marks take their bearing. The boundary tanges $a = a$ once, admits it on the witness, and sets the resulting receipt as the baseline. The mark does not move, and every later structure indexes from it.

This first receipt is the first mark: the initial tick of the apparatus, the calibration baseline $\alpha$ against which later receipts become comparable. It is finite, decided, and held, and it is also mute. What it does not yet have is a successor. A single mark cannot make a geometry; the boundary needs a way to funge one receipt into the next and to register that the second stands after the first. That is the limit the first surface reaches and cannot cross on its own: it can tange a claim into a held mark, but it cannot yet say that one mark comes after another. The next section forms that step---the truth-phase successor that turns isolated receipts into an ordered sequence.

\section[Witnessed Receipt]{The Witnessed Receipt as Claim Plus Finite Decision}

The preceding section fixed the first measuring surface and gave the first mark the form of a receipt. The next distinction concerns what such a receipt contains, because the receipt is the whole difference between a claim that has been measured and a claim that has only been stated. A bare claim does not yet measure anything. It names a possible condition and asks for admission, but it has not passed through a finite decision, and until it has, the boundary holds nothing but the request. The boundary therefore treats a claim as an offered form, not as a fact already in the record: an offer may be accepted or refused, but it is not yet either, and the work of this section is to watch the boundary turn one offer into the first held fact.

Let \(C\) denote a claim. Before the boundary decides it, \(C\) only specifies what must be settled; it fixes the question and leaves the answer open. It may be well formed, meaningful, and even supported by reasons, yet it remains unmeasured until the boundary can return a definite answer in finitely many steps. This is the chapter's sharpest early distinction, and it is worth stating flatly: a claim can be true and still unmeasured. Truth, if the word even applies before there is an instrument to apply it, is a property the claim might carry on its own; measurement is something the boundary must do and then keep. A measured fact needs the claim together with the answer the boundary can hold, and an answer the boundary cannot reach in finitely many steps is, for this apparatus, no answer at all. The boundary does not resent that limit. The limit is what keeps its readings honest, since an instrument that could pronounce on claims it cannot decide would be reporting beliefs and calling them measurements.

The finite decision supplies that answer. Write \(d\) for a terminating finite decision procedure, a rule the boundary can run to its end. Applied to the claim \(C\), it returns one of two values,
\[
  d(C)\in\{\mathsf{yes},\mathsf{no}\}.
\]
This exposed yes-or-no edge is the needle. It is not a physical compass, and it is not a theorem derived after the fact; it is the single classical edge at which the finite decision shows the boundary what answer it may carry. Everything else in the apparatus is deliberately held in suspense, kept open and comparable rather than settled, but at the needle the boundary commits: for this one claim, under this one decision, the answer is yes or it is no, with nothing in between. Much of the book is spent keeping structures from collapsing to a single truth-value too early, since a premature collapse is exactly how a measuring instrument begins to lie to itself; the needle is the one controlled place where exactly one such collapse is allowed, performed on purpose, and paid for.

The important feature of the decision is not optimism about the claim, and not even its correctness. It is termination. The decision has a finite route to a definite result, so the boundary can accept the outcome as something that belongs to its record rather than waiting on a process that may never return. A procedure that might run forever cannot be a decision here, however reasonable its intent, because the boundary would never be able to hold its answer, and an unmeasured claim and a claim whose decision never halts are, from the boundary's side, the same situation. A claim without a terminating decision remains only a proposed form, an offer the boundary has not yet been able to take up.

Termination is not only a yes-or-no fact; it is a quantity. A decision that halts halts after some definite number of steps, and that number is exactly the cost the previous section named: deciding \(C\) costs
\[
  \#\Gamma_C \in \mathbb{N}
\]
steps, the finite count of moves the boundary spends before the needle reads. The decision is therefore the countable thing the apparatus prices. A cheap claim is one the boundary decides in few steps; an expensive claim is one whose decision runs long but still ends; a claim with no terminating decision has no finite \(\#\Gamma_C\) at all, and falls outside what can be measured. The receipt the boundary is about to form will carry, implicitly, not just an answer but the price of reaching it.

The receipt is the held pair produced once the decision halts:
\[
  \rho = \langle C, d(C)\rangle .
\]
The first component says what was asked; the second says how the boundary settled it. In this pair the claim no longer floats as a bare assertion, and the decision no longer disappears as a private act the moment it finishes. The boundary tanges the two into one held object, and tanging is exactly this act of binding a claim to its decision so that neither can be read without the other. What the boundary holds afterward is the finite object later stages may carry, compare, and place in order. Once held, the receipt funges forward: it circulates to later stages as the carried decision, and no stage that receives it has to decide \(C\) again, because the receipt already carries the answer and the price of having reached it. A measurement, in this apparatus, is something decided once and funged many times.

This pairing also separates three things ordinary usage runs together: warrant, decision, and receipt. A warrant gives a reason or a derivation for a claim, an argument that it ought to be accepted. A decision gives a terminating finite procedure that settles the claim with a yes or no, an act the boundary can actually perform. A receipt is the held result after that decision has occurred, the durable trace that the act took place and what it returned. The three can come apart, and seeing how is the point: a claim may have a strong warrant and no terminating decision, in which case the apparatus cannot measure it however convincing the argument; a claim may be decided and its result discarded, in which case there is no receipt and so, for the boundary, no measurement at all. A warrant may explain why a claim should be accepted, but measurement here needs a stronger and humbler fact, custody rather than persuasion: the claim met a finite decision, the decision halted, and the result was kept where a later stage can find it.

That custody is what makes a receipt a unit rather than a passing occurrence. Because the receipt holds both a claim and the decision that settled it, two receipts can be set side by side and compared as objects without either claim being reopened, and the comparison reads what was decided rather than the arguments behind it. This is why the held pair, and not the bare answer, is the thing the rest of the chapter carries: an answer alone could be reproduced but not located, while a receipt can be tanged once and then funged from stage to stage, placed, found again, and ranked against another. A stage that receives a funged receipt inherits a settled fact rather than an open question, so the work of deciding is done once while the work of comparing can be done as often as the construction needs; the receipt is what lets the apparatus pay the decision's price a single time and draw on it indefinitely. Order, when it arrives, will be an order on receipts, which is exactly why the receipts must be tanged and held before any order can be built on them. The first example shows the smallest case of all.

The first example is still reflexive equality, the cheapest claim that carries a genuine decision. For the calibration claim \(a=a\), the boundary decides by presenting the same object on both sides, so the decision halts in a step and returns \(\mathsf{yes}\), and the first receipt takes the form
\[
  \rho_\alpha = \langle a=a,\mathsf{yes}\rangle .
\]
This receipt is the baseline \(\alpha\). It does not establish a world beyond the boundary, and it does not claim that \(a=a\) is true in any sense larger than the one the boundary just checked. It establishes the one thing the chapter has been working toward: that the boundary can tange a decided claim into a held receipt, and so can hold a fact rather than merely entertain a claim. Everything the book builds after this point rests on receipts of this kind and never on claims that were only argued for, which is why the smallness of the baseline matters as much as its existence: it is the cheapest thing the apparatus can stand on, and the apparatus will stand on nothing it has not paid for in the same way.

Once tanged, the baseline can serve as an origin for later construction, and this is where the receipt begins to funge. The claim \(a=a\) need not be introduced again as an unsettled proposal each time the chapter uses it; the boundary funges \(\rho_\alpha\) forward, and what circulates is the decided claim, not the open question. The receipt carries that decision as a stable mark, and its stability is modest and exact. It does not say that \(a=a\) is necessary, or eternal, or true beyond the boundary. It says only this: that this claim, under this boundary, has a finite yes-decision attached to it, and that the decision and its price have been kept. Modest as that is, it is the first thing in the book the apparatus can stand on.

That modest stability is enough for the next step, and not one step more. A single held receipt gives the chapter an origin, but it gives no succession; the baseline can say that it is, and cannot yet say what comes after it. Geometry needs a way for one receipt to stand after another without losing the finite decision that made each receipt measurable, so that order is added to custody and not bought by giving custody up. That is the exact thing a single receipt cannot supply from inside itself, and the exact thing the next section supplies: a successor structure that is not neutral arithmetic but an ordered phase of receipts beginning from the baseline, in which the second receipt knows it stands after the first and still carries its own decision intact.


\section{Truth-Phase Successor Rather Than Neutral Arithmetic}

The baseline receipt gives the chapter one held mark, and one mark is not yet a geometry. A geometry needs a way for a later receipt to stand after an earlier one while both stay attached to their finite decisions, so that order is gained without custody being given up. The truth-phase successor supplies exactly that, and it is the first place the book insists on a distinction it will defend the whole way through: the successor is not neutral arithmetic imported from elsewhere. No number line is borrowed, no counting is assumed to exist before the boundary builds it, and no position is read off a scale that was lying around in advance. What the boundary builds instead is a finite receipt chain that begins from a nonempty origin and adds one receipt at a time, and the number that emerges from it is a phase of that chain, an index of how far the record has been carried, not a quantity living outside the record and applied to it.

The difference matters because an imported number would smuggle in exactly the thing the apparatus is trying to earn. To borrow the natural numbers whole would be to assume, before any measurement, that a settled sequence already exists for the boundary to consult; but the boundary's entire discipline is that it consults nothing it has not decided and tanged. So it builds its succession the only way it is allowed to, out of receipts it already holds, and the order it gets is therefore its own, answerable at every step to a finite decision rather than to a scale taken on faith.

The origin is nonempty because the boundary cannot order a void. There is nothing to put before or after in an empty record, so succession has to begin from something already held, and the natural something is a receipt the boundary has already decided and tanged, such as the fixed reference \(\rho_\alpha\) from the preceding section. The first position therefore carries a decided claim. It is not the absence of value, and not a zero waiting to be filled; it is the initial phase of the record, the place from which every later position will count its distance. Beginning from a decided receipt rather than from emptiness is not a convenience to be optimized away later. It is what lets every later number stay a phase of truth rather than a mark on an imported scale, and it is why the chapter spent a whole section tanging the baseline before it allowed itself a second mark.

That word phase can be made exact, and once it is exact it is countable. Write \(n\) for the number of successor steps taken from the origin. The position reached after \(n\) of them is the \(n\)-th truth-phase,
\[
  n \in \mathbb{N},
\]
and \(n\) is the only number this section produces. It does not measure a magnitude and it does not point at a place on a line; it counts how many decided receipts have been tanged into the chain since the origin. The phase is the count, and the count is finite at every position, because each step adds exactly one receipt and no step adds anything the boundary has not already decided. A number, in this book, is first of all a tally of held decisions, and only much later, once enough has been built, anything that resembles a magnitude.

The successor step adds a new receipt to a previous state. If \(s\) names an existing state and \(\rho\) names a fresh receipt, the next state may be read schematically as
\[
  \mathsf{succ}(\rho,s).
\]
This notation does not assert ordinary addition, and the difference is again the section's whole point. \(\mathsf{succ}(\rho,s)\) does not combine two numbers into a larger one; it tanges the current receipt together with the earlier state into a single later state, so that the new position holds its own decision and carries, intact, the chain it grew from. Nothing is summed and nothing is discarded. Repeating the step produces a finite nested chain: the origin, the successor of the origin, the successor of that successor, and so on, each position answerable to the receipt at its outer edge and to the earlier state folded inside it. The chain funges forward as it grows. A later state inherits every decision tanged into the states before it, and inherits them without re-deciding any of them, exactly as a single receipt was decided once and then funged many times; succession here is just custody extended one receipt at a time, and the number \(n\) is the length of that custody.

Because the chain is built and not borrowed, it carries something an imported number could not: a memory of how each position was reached. A neutral integer three says nothing about what made it three; this chain's third phase still holds the three decided receipts that brought it there, in order, each available to be read again. That memory is not decoration. It is what the order test is about to use.

The chain gives succession, but succession alone does not yet compare two states. A record that knows how to grow does not, merely by growing, know how to say which of two growths stands first, and the boundary may not answer that question by stepping outside itself to measure. It needs a finite order test that stays inside the same discipline that built the chain: a rule that may read decisions and follow the chain inward, and may do nothing else. The base rules are direct, and they fall out of what the chain is. An origin precedes or equals any state, because the origin is where every chain starts and nothing can stand before the place the others were tanged from. A successor does not precede the origin, because nothing added after the origin can come before it. The only case left is the one that carries the section's real content: two states that are both successors, whose order must be read off the decisions they carry rather than from any position assigned in advance.

Let \(p \preceq q\) name this local comparison. When the two inputs are successor states, their outer decisions determine the next step:
\[
\begin{array}{c|c|c}
d(\rho_1) & d(\rho_2) &
\mathsf{succ}(\rho_1,p)\preceq\mathsf{succ}(\rho_2,q) \\
\hline
\mathsf{yes} & \mathsf{yes} & p\preceq q \\
\mathsf{yes} & \mathsf{no}  & \mathsf{no} \\
\mathsf{no}  & \mathsf{yes} & \mathsf{yes} \\
\mathsf{no}  & \mathsf{no}  & \neg(p\preceq q)
\end{array}
\]
The table is the whole local orientation rule, and it rewards reading slowly, because it is the first appearance of a pattern the book returns to constantly. Two affirmative receipts pass the comparison inward unchanged. If both outer decisions read \(\mathsf{yes}\), the boundary asks the same question of the earlier states, \(p \preceq q\), and the orientation is preserved as the comparison moves one step toward the origin. This is the covariant case: the inner comparison runs the same way as the outer one, and an affirmative edge is, in effect, a decision that the order found at this position should carry on unchanged. Two negative receipts reverse it. If both outer decisions read \(\mathsf{no}\), the boundary still asks about the earlier states, but it now asks the opposite question, \(\neg(p \preceq q)\), and the orientation is flipped. This is the contravariant case: carrying a negative decision through the comparison turns it around, so that the same inner states are read in the opposite sense.

The two mixed cases never go inward at all. An affirmative receipt set against a negative one fails immediately, a negative set against an affirmative one succeeds immediately, and in both the chain below is never consulted, because the outer decisions already disagree and the disagreement is itself the answer. The boundary does not need to look further to know that a position carrying \(\mathsf{yes}\) and a position carrying \(\mathsf{no}\) are oriented against each other; the edge has decided the matter, and a finite test that can stop early stops.

It is worth being plain about what the covariant and contravariant cases are, because they will reappear under heavier names later. They are the record's memory of which decisions ran with the order and which ran against it. The boundary refuses to pretend that affirmation and negation point the same way, so when it carries a comparison inward it carries with it whether the current edge agreed or disagreed with the direction of the order. Two agreements compound and the order is preserved; two disagreements compound and the order reverses; one of each is a contradiction at the edge and resolves there. This is not a quirk of the notation. It is the smallest honest bookkeeping for a record that decides each step rather than reading positions off a line, and the whole later apparatus of covariant and contravariant comparison is this same table, kept and reused.

Nothing in the rule subtracts positions, measures a numerical gap, or appeals to an outside number line. It reads the carried decisions at the current edge of the two chains and then either stops, when the edge decides the matter, or continues inward, when the edge passes the question along. The order it produces is finite and receipt-bound, an order claim made by the same boundary discipline that produced the receipts themselves, and it is in this sense that the comparison is honest: it never imports a fact the boundary did not already tange, and it never reaches a verdict the held decisions did not already contain.

At this point the chapter has an origin, a successor step that tanges one receipt onto the last, and a finite order comparison that funges decisions inward without ever leaving the record. It can say that one receipt chain stands before another, and say so for chains of any finite length, using nothing but the decisions already held. What it still cannot do is carry a distinct measured value. The phase counts position, but position is not yet magnitude; a chain that knows its own order does not yet know what travels along it, and a number that is only a tally of decided receipts is not yet a number that can vary while the order stays fixed. That gap is deliberate, and it is the exact thing the next section opens: a vessel, a local place where a carried value can travel along the ordered receipt chain without being confused with the chain that carries it.


\section{Vessel versus Carried Value}

Succession orders one receipt after another, but ordering alone fixes no place at which a value can sit. It says that a receipt follows a receipt; it does not say where the present sign is exposed, how that sign differs from the history already gathered, or how the next update keeps the two from collapsing into each other. The vessel supplies the missing place. It introduces no new number line and no hidden magnitude. It gives the apparatus a local holder in which a present label and a carried mark stay distinct while belonging to one record.

A vessel does two things, and only two. It names the local place under inspection, and it keeps the current label apart from the carried mark. The current label is the apparatus-facing sign exposed at the present edge of the vessel. It is affirmative or negative --- the same truth-phase sign the most recent successor step set down --- and it is the sign showing now, at this place, and nothing more. It is not the receipt chain, and it carries no history of its own. What the vessel holds, behind that exposed edge, is a tanged mark: a record fixed in place by an act of measurement and kept available for later inspection.

The carried mark plays the other role. It is the receipt chain already gathered for this vessel, the accumulation that earlier succession produced. It is not an integer, not a magnitude, not a hidden value on some continuum waiting to be read off. Nor is it reached by an external index: the apparatus does not address the carried record by a coordinate but recovers it only by undoing wrappers, peeling the enclosure one layer at a time. It is not a pointer to a value kept elsewhere, because there is no elsewhere; the mark is the value, held in place, and the only route to the history it encodes runs through the wrappers themselves. It is a finite mark, carried forward from the baseline receipt through its successors, one wrapper at a step. The separation stays strict --- the label reports what the vessel presently shows, and the carried mark reports what the vessel has already brought along. One is the face; the other is the freight.

Holding the label apart from the carried mark is what makes the vessel more than a receipt under another name. A reading can change at the edge without altering what is held beneath it, and what is held can grow by one wrapper without changing how the edge is read. The two components vary independently, yet they belong to one holder: a single local place that presents a face and carries a freight at once. Keeping them distinct is exactly what lets the apparatus track a present sign and an accumulated history in the same record without confusing the one for the other.

Write the vessel as the pair of its two parts, the exposed label together with the carried mark, and write the update that advances it:
\[ v=(\ell,m), \qquad v'=\bigl(\ell',\,\mathrm{wrap}(m)\bigr), \qquad \ell\in\{\oplus,\ominus\}. \]
At the origin the vessel carries the baseline receipt, with no wrapper yet placed around it --- the ground case from which every later carried mark is built. At a successor the update takes the mark already held, $m$, and places exactly one new wrapper around it, $\mathrm{wrap}(m)$, while the label is re-read at the new edge as $\ell'$. The step touches the carried component alone. It is local in the strict sense: no subtraction, no distance, no external coordinate, and no reach to any place but this one. Exactly one wrapper is added, never several at once and never none, so each update advances the record by a single definite layer.

That single wrapper moves the boundary by exactly one layer. The new state does not merely point past the old one. It holds the old mark inside the new mark, so the edge at which the apparatus reads the vessel has shifted one step outward along the receipt chain while the earlier record stays enclosed within it. The carried marks therefore form a nested sequence, each enclosure holding all the earlier ones inside it. Because each step adds a layer and removes none, the depth of the enclosure only increases: the carried record grows along the chain, and no update returns the vessel to a shallower state. This is the whole displacement claim: the exposed edge of the carried record advances, and nothing of what was carried is lost. The value funges forward --- it passes from the vessel as it stood to the vessel as it now stands, carried across the update rather than recomputed at it.

The distinction is between the holder and what is held. A vessel is a standing site; the value is what funges through it. The two never merge. A holder keeps its identity across an update while the carried mark moves on, and the carried mark keeps its content while the holder it occupies stays fixed in place. The carrier is not the carried.

A single vessel now provides a local holder, a present label, a carried mark, and an update that advances one wrapper at a time. The apparatus can say what this place shows and what this place carries, and it can funge the carried record forward without confusing the showing with the carrying. The next difficulty appears the moment one vessel is not enough. To compare several such places --- to count them, order them, or track one against another --- the apparatus must first tell one vessel from another by finite means, before any of that comparison can begin.


\section[Convergence Promise]{Convergence Promise and Residue as the First Hint of a Limit}

A single vessel gives the apparatus one place to read, but a record worth pressing toward a limit needs more than one. The moment several vessels stand together, a new discipline is forced: before the apparatus can count them, order them, or follow any across time, it must tell which place carries which receipt. The test it needs is local and modest. It does not settle identity in the abstract, and it never asks what a vessel is in itself; it asks only whether the candidate now under inspection is this prepared place and not that one --- a bounded yes-or-no read straight off the marks the vessels already carry. Distinguishability of this finite kind is operational, not metaphysical: the apparatus requires no global account of sameness, only a discriminator sharp enough to keep two prepared places from being confused at the moment of reading. And it is the precondition for everything the section now builds, because a narrowing has no meaning until the apparatus can say which records are the ones narrowing, and can find each of them again when it returns.

With places held apart, the apparatus lays down a tally. The tally imports no ready-made number line; it begins at an origin and advances by a single step, and each step tanges one more receipt into a finite chain. Its order is read, not assumed: one counted record precedes another when the decisions it carries and the truth-phase already built in the first mark place it earlier --- never because an infinite arithmetic background has quietly appeared to rank the records for it. The difference matters. An imported number line would let the apparatus consult an order it never decided, ranking its records against positions it did not tange; the built tally consults only what it has already laid down. It is therefore a number with a memory: it records not merely how far the count has gone but the order of the steps that took it there, and that memory is exactly what the later comparisons read when they ask which record came first.

Counting becomes an active local act rather than a lookup against a fixed scale. A count holds three things at once: the vessel under inspection, the tally reached so far, and the step by which the apparatus moves to the next prepared place. To these an admissibility condition adds the single thing finitude demands: a bound. Admissibility does not prove that every conceivable search terminates, and it does not shield the apparatus against every divergence that might arise; it makes no global promise of any kind. It says only that the proposed count carries a local condition, read from earlier receipts, under which the apparatus may stop inside the finite ladder it has already prepared. The bound is part of the record rather than a hope held about it from outside --- a certificate the count carries with it, so that stopping is a decision the record already contains and not an intervention imposed on it.

Comparison at this scale also needs proportions, and the apparatus builds them as finitely as everything else it owns. A finite proportion is a pair of tallies read as numerator and denominator, with zero kept as its own base case rather than treated as a forbidden division. It is not a field of rational numbers, and the distinction is deliberate: the apparatus declines division, completion, and continuous magnitude, yet it keeps the shape of a ratio --- a comparison of one bounded share against another, weighed without ever forming a quotient. Nothing continuous enters the chapter through this door; the proportion is a way of ranking prepared shares, not a number on a line waiting to be evaluated. The order the apparatus reads on two such proportions is fixed by a finite rule:

\[
\begin{array}{c|c}
\text{case} & r \preceq_{\mathrm{frac}} s \\
\hline
r=0,\ s\neq 0 & \mathsf{yes} \\
r\neq 0,\ s=0 & \mathsf{no} \\
\text{same decision} & \text{compare both tallies forward} \\
\text{opposed decision} & \text{reverse the carried comparison}
\end{array}
\]

The rule is nothing more than a finite ordering, and it carries the same orientation the chapter has used since the first successor. It tells the apparatus how to rank two prepared proportions once it has read the decisions they carry, and no sooner. When the two carry the same decision, numerator and denominator move with the order already recorded and the comparison runs forward; this is the covariant case, the ranking that preserves the carried order. When they carry opposed decisions the comparison reverses, because the receipts point against one another and the order they imply is inverted; this is the contravariant case. The base cases anchor the rest: a zero numerator against a nonzero denominator ranks below, a nonzero numerator against a zero denominator ranks above, and at no point does the rule appeal to a magnitude the proportions do not already carry. The order is read out of the receipts, exactly as the tally's order was.

A list of such proportions is useful only if the apparatus can find its way back to any one of them, so indexing comes next. An index carries an origin, a step, and a bound, and it makes no claim to infinite countability; it is a finite address book, not a census of all possible places. Its single task is to keep each prepared value fastened to a finite place, so that a later reader can ask for the next value without losing the order of the list. An encoded sequence is exactly that ordered list together with the rule that makes it readable --- a finite arrangement the apparatus can traverse one prepared value at a time. What the index buys is precisely the ability of a finite reader to stand in for an inexhaustible one: the apparatus never holds the whole of a sequence at once, only a place in it and a step to the next, and that is enough to read the order without ever completing it.

Tracking reads the encoded sequence under a named bound. It is not yet the extraction of a limit, and it does not pretend to be one. The tracker holds four things together --- the sequence, the finite proportion currently in view, the boundary accumulated so far, and the step that advances the read --- and only when those four cohere does the apparatus permit itself a convergence promise. Coherence here is concrete: each prepared stage must sit no further out than the stage before it, so that the residues read along the way do not grow. Writing $r_k$ for the residue standing at prepared stage $k$, the promise is the non-increasing chain
\[ \cdots \;\preceq_{\mathrm{frac}}\; r_{k+1} \;\preceq_{\mathrm{frac}}\; r_k \;\preceq_{\mathrm{frac}}\; \cdots \;\preceq_{\mathrm{frac}}\; r_0, \qquad r_k \neq 0 . \]
Each $r_k$ is a finite proportion ranked by the rule just fixed; the chain narrows, while the standing $r_k \neq 0$ records that the narrowing is pressed and never crossed. The narrowing tanges the ordered record toward a tighter interface, stage by stage, and at every stage the promise can be checked against a representative without appeal to anything outside the finite. What the promise is not must be said plainly, because the temptation to overread it is strong: it is no continuum limit, no norm, no metric, no gauge, and no completed analytic result. The chain could narrow through every prepared stage and still hand the apparatus nothing at the end, and the apparatus never claims otherwise. It is a bounded claim about the records now in hand, and it stays a claim about records.

The promise rests on nothing the chapter has not already built. The bound that makes a count stop is the same discipline that lets a narrowing be checked at a finite stage: in both cases the apparatus reads a local condition off the receipts and acts on it without consulting anything global. The convergence promise is admissibility wearing a second face --- a finite certificate, carried by the record, that this stage may be examined and the next deferred. Nothing in it reaches past the finite ladder; it only reads further along rungs already in place.

Residue names what remains when the record can be funged no further. At each prepared stage the narrowing leaves something over: a finite proportion still standing above the base case, paired with a representative the apparatus can actually exhibit in place of the limit it does not hold. Written as a pair,
\[ \mathrm{res}_k \;=\; (\,r_k,\ \sigma_k\,), \qquad r_k \neq 0 , \]
with $r_k$ the leftover interface and $\sigma_k$ the bounded sample that stands for it. The residue does not finish the measurement; it is the part that funges forward when nothing more can be funged at this stage --- the finite gap carried across to the next act. Two things are true of it at once, and both matter: it keeps the finite promise visible, so the narrowing already done is not lost, and it marks the gap the next act must still read, so the work is honestly recorded as unfinished. This is the first hint of a limit the chapter permits itself --- not the limit, but the standing evidence that one is being approached and has not been reached. The residue is the interface between what has been pressed and what remains, and it stays finite enough to be read again.

The last object the chapter assembles is a clocked observation. It names an initial condition and a response, together with the finite order rules that compare condition with condition, response with response, and response against condition. The first two rank like with like; the third is the chapter's earliest comparison of one kind of mark against another, and it is the seed the later work grows from. The observation stays deliberately modest: it is not a completed measurement, not a universal quantity, and not a continuum result. It is the first clock-ready record built from the whole apparatus this chapter has prepared --- vessels and tallies, counts and proportions, indexes and the encoded sequence, tracking, the convergence promise, and the residue. With it the first mark has become a record that can point past itself while staying entirely inside the finite. That is the chapter's quiet result: the promise of a limit is itself a finite object, and the residue is what keeps the promise honest. A limit named is not a limit reached, and the difference between them is not an embarrassment to be completed away but the very thing the apparatus carries forward. The bridge now asks what can observe such a record.


\section*{Bridge: Residue Left After the First Mark}

Chapter 1 ends with a residue: a finite observation interface that holds an initial condition beside its response. Everything the first mark could produce on its own is gathered there. The whole apparatus of the chapter --- vessels and tallies, proportions and indexes, the encoded sequence and the convergence promise --- works to tange a condition beside a response and to fix that pair in place as a bounded record, ordered by the same finite rules the chapter built one at a time. That record is real, it is finite, and it is the most the construction can make for itself. But it carries one incapacity no further step inside the chapter can remove: it cannot observe itself. It holds a prior state beside a later state, and it even holds the rule that would compare the two, yet holding a comparison is not performing it. The record has no means to set its earlier state against its later state and report that a transition occurred. Tanging a record is not the same as reading it, and the first chapter only ever learned to tange. It built a record that points toward a transition and stops short of pronouncing one, because pronouncing requires a standpoint the construction never took.

The missing thing is an observer, and Chapter 1 has not supplied one --- not by oversight, but by the shape of what it built. The residue is the object to be observed, not the act of observing it, and the gap between the two does not close by adding one more vessel or one more rung to the tally. Every object the chapter assembled is something the construction makes for itself: a forward order, laid down step by step, read off receipts it has already decided. That forward order is the construction's own clock --- the order in which it lays its marks down. Reading the record back, from an earlier state toward a later one as an observer would, asks for the order run the other way, from the side the construction never had to occupy because it was always moving forward. An observer stands in exactly that complementary relation. It does not extend the record from the inside; it reads the record from the side the construction cannot reach. What the chapter funges forward, then, is a finished interface and a single open question: if a bounded observation exists, what observes it?

That question sets a strict condition on the observer the next chapter must introduce, and the condition is sharp on both sides at once. On one side, the observer has to record a difference between an earlier state and a later one. A reader that cannot tell before from after registers nothing, and the residue stays inert beneath it. The forward order the construction laid down answers only what comes next; an observer must answer a different question --- whether the later state departs from the earlier one at all --- and that question is the observer's to ask, not the construction's. The difference it records stays as finite as everything the chapter has built. The observer does not measure how far the later state departs from the earlier one, as though change were a magnitude on a scale; it registers only that they depart, a bounded distinction read off the record and nothing more. To observe is to set the prior state against the later state and find them not the same.

On the other side, the observer must not fix the truth too early. The record the first mark leaves is built on a truth-phase carried forward rather than pronounced, held in suspension precisely so the order can be read without the verdict being forced. An observer that resolves that truth before the phase is tracked collapses the very distinction it was meant to read. Resolve too soon, and there is nothing left to observe but a verdict the observer itself has imposed; resolve nothing, and the observation never happens at all. The suspension is not indecision but capacity: a phase held open can still be read more than one way, while a phase collapsed too early admits only the single reading the observer forced upon it. The observer must therefore hold a narrow line between two failures --- sensitive enough to register that a state has changed, restrained enough not to decide, in the act of registering, what the change finally means. The needle from the first receipt remains, at the close of the chapter, only an exposed decision edge: not yet a sign, not yet a charge, not yet a completed reading. It is poised over the record and not pressed into it, and it waits to be read without being forced.

This is the whole of what the first chapter hands forward. The first mark is made; the record is tanged and bounded; the residue funges forward and waits. What it waits for is an observer that can read the construction's forward order from the complementary side --- recording that an earlier state and a later state differ, without collapsing the phase in which the difference is carried. No part of the apparatus built so far can be that observer, because every part was made to lay the order down and not to read it back. The reader the residue waits for is not inside the record it would read; it must come from the next chapter, as the complement of the order this one laid down. The first mark is made and the residue waits; the problem of the observer begins.

\section*{Coda: The First Mark in Review}

Chapter 1 carried a single claim across a finite measuring surface and left a residue. The route runs through a fixed order of plain steps, and this coda walks that order back --- not to summarize it for convenience, but to show that it was forced. Each object in the chain arose because the step before it exposed something it could not handle, and none was imported from outside the finite surface. Read forward, the chapter looks like a list of definitions; read backward, it is a list of answers, each one the least the apparatus could add and still keep going. The review below follows the forward order and marks, at each step, the lack that made the next step necessary.

The route begins at a baseline: reflexive equality, the assertion that an entity is identical to itself. This serves as the chapter's calibration point, the fixed place from which every later mark takes its distance. It holds as a plain fact, and that is exactly its limit --- a fact asserted is not a measurement made, because nothing has yet been decided about it or carried away from it. The baseline tanges nothing; it stands still and waits to be marked. What it supplies is an origin, the one position the apparatus may treat as given, so that later receipts have somewhere to be counted from. A measurement that began nowhere could not report distance, and the baseline is the chapter's answer to that lack; its own insufficiency --- a place to stand, but nothing yet standing there --- is what calls for the first mark.

The first mark tanges that baseline into a receipt. A receipt is a claim together with a finite decision about it --- an explicit carried witness rather than raw, unattached truth. At the exposed yes-or-no edge of that decision the needle names the answer the apparatus can carry, and the count of marks gathered there, $\#\Gamma$, is the only number the receipt reports. The apparatus accepts a carried witness and not a bare assertion for a reason the rest of the chapter depends on: only a witness that has been decided and held can be inspected, challenged, and funged onward. Raw truth, however certain, cannot be passed from hand to hand until it is fixed, and what cannot be passed cannot be measured. The receipt is the chapter's answer to that requirement --- it turns a standing fact into a portable record --- and it raises the next lack at once, because a single receipt has no place in any order until something arranges it.

The truth-phase successor places the receipt in order. Written $\mathrm{succ}(\rho,s)$, it carries a receipt $\rho$ forward under a sign $s$, and the number it produces, $n \in \mathbb{N}$, counts not quantity in the ordinary arithmetic sense but how many discrete phases separate the present receipt from the baseline. It reports where the route stands, not how large a magnitude it has reached. The order it lays down is read off the receipts themselves, never imported from a number line the apparatus did not construct, and it carries an orientation: receipts that agree run with the order, the covariant case, and receipts that oppose it run against, the contravariant case. The successor is the chapter's answer to the receipt's placelessness. But an order of receipts still says nothing about where each one is held, and that silence is what the vessel is built to break.

A vessel gives the mark a place to stand. Written as a pair $v=(\ell,m)$, it holds an exposed label $\ell$ beside a carried mark $m$ --- a localized arena that ties the receipt to one property space while staying strictly distinct from the value it carries. The same carried value can rest in different vessels, and one vessel can receive different values; holder and held never merge. An update wraps the carried mark, $m \mapsto \mathrm{wrap}(m)$, advancing the record one layer without disturbing the label read at the edge, so that the carrier is not the carried. The vessel is the chapter's answer to where a receipt is held. Yet a single vessel, however neatly it separates face from freight, cannot be set against another until the apparatus can tell the two apart, and that is the next lack the route must close.

Distinction separates one vessel from another. The route needs no global map of every vessel and no account of identity in general; it needs only a local difference --- a distinguished pair --- sharp enough that two marks are never silently merged at the moment of reading. The discrimination is operational, not metaphysical: it asks only whether the place now under inspection is this one and not that one, a bounded yes-or-no read off the marks the vessels already carry. Distinction is the chapter's answer to the bare requirement of comparison, and it is the precondition for everything quantitative that follows, because a narrowing has no meaning until the apparatus can say which records are the ones being narrowed and can find each of them again.

With places held apart, a count gathers the distinguished marks. The count stays within a local bound --- an admissibility condition read from earlier receipts --- so the enumeration never outruns the finite surface that carries it; the bound is part of the record rather than a hope held about it, and it lets the apparatus stop inside the ladder it has prepared. To weigh one bounded share against another the count forms finite proportions, pairs of tallies ranked by a finite rule with zero kept as its own base case, which hold the shape of a ratio while declining division, completion, and continuous magnitude. From the count an ordered sequence follows: an arrangement with an origin, a step, and a threshold, indexed so that each prepared value stays fastened to a finite place and a later reader can ask for the next without losing the order. None of this imports an infinite background; the apparatus never holds the whole of a sequence at once, only a place in it and a step to the next, which is enough to read the order without completing it.

Tracking reads that sequence under a named bound and, when its parts cohere, states a convergence promise: the residues read at successive stages do not grow, so the record presses inward toward a tighter interface. The promise is the non-increasing chain $r_{k+1} \preceq_{\mathrm{frac}} r_k$ holding back to $r_0$, with each $r_k \neq 0$ --- a narrowing pressed and never crossed. The narrowing tanges the ordered record toward that interface, stage by stage, and the sequence leaves a residue when it can be funged no further. The residue is a finite interface that pairs the convergence promise with a representative the apparatus can exhibit, $\mathrm{res}_k = (r_k,\sigma_k)$, so that a bounded sample stands in for the limit the apparatus does not hold. It is not a completed numeric result, an analytic convergence, a norm, a metric, or a gauge; it answers one bounded question about its own limits and funges forward to whatever can read it. The residue is where the chapter's construction runs out of moves.

Set out this way, the order shows itself to be a dependency chain rather than a convenient sequence. The apparatus cannot place a mark before it has a baseline to mark from; it cannot order receipts before there are receipts to order; it cannot hold a value before there is a vessel to hold it; it cannot compare vessels before it can tell them apart; it cannot count before it can compare; it cannot press a narrowing before it can count; and it cannot leave a residue before a narrowing has been pressed. Every step strictly presupposes the one before it, and no step reaches past the finite surface for outside help. That is what makes the route minimal --- not that it is the shortest path anyone might draw, but that nothing in it can be dropped or reordered without dissolving a step that depends on it. Each object is the least the apparatus could add and still carry the claim one stage further. The chapter is not a list of conveniences the construction happened to adopt; it is the smallest finite scaffolding under which a single claim can be marked, carried, ordered, and pressed to the point of leaving a residue.

In plain form, the chapter's progression is mark $\rightarrow$ receipt $\rightarrow$ vessel $\rightarrow$ distinguished pair $\rightarrow$ count $\rightarrow$ sequence $\rightarrow$ residue.

It is worth recalling what all of this was in service of. The chapter set out to carry a single claim across a finite measuring surface, and the residue is that claim at the end of its carry --- decided, held, ordered, counted, and pressed, now standing as a finite interface rather than as the bare assertion it began as. Nothing was added to the claim except the scaffolding required to move it, and nothing continuous was smuggled in to complete it. What the apparatus holds at the close is the same claim it started with, changed only in that it can now be inspected and funged onward, and halted only because the one thing it still lacks is a reader. The whole finite route was the cost of turning an assertion into something an observer could one day read.

The residue is where Chapter 1 stops, and the reason it stops is exact rather than incidental. The interface it leaves cannot observe itself: it holds an initial condition and a structural response side by side, and it even holds the rules that would compare them, but it has no means to perform that comparison or to report whether a transition occurred. The finite surface can do a great deal --- place a mark, tange it into a receipt, order the receipt by its phase, hold it in a vessel, distinguish one mark from another, count, form proportions, index, track, and deposit a residue --- and at the end of all of it the surface stays blind to its own result. Every step was an act the construction performed on itself, in the forward order it laid down; not one of them was a reading of that order from outside. That blindness is not a defect to be patched but the chapter's honest boundary: a finite construction can produce a record and cannot, from within, become the observer of it. The residue funges forward intact, and the question it carries --- what observes the residue --- is one the apparatus cannot answer in its own terms. Chapter 1 closes with the finite interface complete in hand and unread: the first mark is made, and the eye that would read it is still absent.
